\section{Training, policy costs, and deployment in the investment economy}\label{sec:study69}
The economic primitives, feasible set and continuous innovation laws are those of Sections~\ref{sec:accuracy62} and \ref{sec:query63}. The new experiment changes the fitted proposal, verification workload and evaluation precision, not the economy. Every callable controller acquires the state downward and returns only certificate-authorized feasible actions. The hypothesis that matters for deployment is a joint statement about achieved policy cost and the complete work of returning that policy.

The preceding comparison is retained in full in the complete development companion, with its original records and inference families. It found adverse cold-start ReLU timings, no ReLU reuse win, only a late and task-specific quadratic crossing, and unresolved cost differences. The original large four-date vector calculation and the corrected cumulative-array counterexample remain visible. New experiments do not overwrite those findings or reinterpret repeated observations as new samples.

\subsection{A prospectively fixed comparison}
The scalar tasks have $(d,T)=(2,2)$ and $(8,3)$, with original-optimum loss targets $1/128$ and $1/64$. They use twenty-four-bit acquired states and a thirty-two-bit downward action lattice. The certificate-aware ReLU and Bernstein readouts, the label-loss ReLU and the quadratic proposal receive identical purchased training information. Two fixed fitting seeds and four separate process repetitions are declared in advance. Adaptive and upper-chord minimization remain prediction-free controls. The earlier linear, spline, nearest-action, derivative, bisection and unrestricted ReLU comparisons retain their original place in the historical catalogue; they are not described as newly rerun arms.

Independent validation purchases 256 fresh contexts at every nonterminal date. The sixty-four-block reuse design fixes new workload streams before deployment, pays for every model load and identity check, and includes the declared stale-identifier fallback. The two-control comparison uses the same coupled-capacity economy and two independent innovations, with adaptive and uniform triangular refinement and two trained two-output proposals. Four-date neural proposals use the explicitly declared three-date transfer.

All planned parent-process receipts, finite optimizer statuses, warning messages, timeouts and return failures belong to the catalogue. Counts of services are not probabilities of successful learning. Model identities, data, exact optimization certificates and output traces bind each policy to its own construction.

\input{tables/findings69}

\subsection{Original policy costs at a declared precision}
The inference experiment compares adaptive search, the two newly certificate-trained readouts, and the original frozen ReLU callable controller from the preceding study. The last arm loads its original model rather than retraining it. There are 1,024 new common-path rows per task under the original continuous initial and innovation laws. Forty-eight-bit continuous bins enclose the sampled paths. The finer evaluator tolerance is fixed at one thirty-second of the policy's declared target.

The centered observations in Theorem~\ref{thm:centered69} have the original policy-cost difference as their expectation. Their smaller support follows from verified local regret, while their numerical enclosure requires additional original Bellman evaluations. The ordinary cumulative cash-cost estimator is computed on exactly the same new paths. The comparison therefore exposes the gain in resolving power and the additional evaluation work, rather than confusing a regret bound with a measured economic difference.

\input{tables/centered69}

The two-sided design margin is $1/256$ in normalized economic-cost units. It was fixed before the new contrasts, with numerical and sampling half-width allocations of $1/512$ each. The supplementary precision table reports whether each allocation was achieved. Equivalence is asserted only for intervals contained in the declared margin, and noninferiority only when the relevant upper endpoint satisfies it. These findings do not identify a strict policy gain and do not calibrate the normalized margin as money. A negative or positive empirical mean alone licenses neither adoption nor rejection.

The numerical ledger distinguishes the conditional Bellman interval, the state-box expansion, unresolved acquisition branches and final endpoint arithmetic. The state-box expansion combines continuous-bin uncertainty with propagated state arithmetic; those two components are not separately identified in the implementation. Intersections with valid local-regret bounds can reduce the final width, so raw component widths before intersection need not add to the printed final interval. Ambiguous acquisition branches are conservatively retained, not omitted from the sample. Combining the separately allocated validation, centered-cost and cash-cost families requires their union bound.

\subsection{Achieved training and complete cold-start work}
The supplementary readout table reports the achieved convex objective and exact box dual gap at the returned quantized weights. The validation table compares the work-loss account with actual full-action recovery queries on predictor-independent contexts. A validated root-query upper bound is not a bound on process seconds, and a valid policy return is not evidence that a hidden-layer optimizer converges reliably.

\input{tables/cold69}

The complete cold-start clock begins before process launch and includes imports, purchased training information, fitting and exact readout certification, prediction, all Bellman descendants, path enclosure, serialization, durable output and shutdown. Parent receipt bookkeeping is additionally charged in the cohort total. Machine-independent query counts, rational recovery, training queries and largest live batches remain separate fields. The timing range includes every declared repetition. Available frequency and governor information is recorded, but frequency is not claimed experimentally locked; close clock differences are consequently descriptive rather than identified hardware-invariant gains.

\subsection{Reuse after paying for the model}
\input{tables/reuse69}
A final complete-process win and an early persistent prefix crossing are different events. The former includes shutdown and the final output account; the latter requires a crossing by block forty-eight that survives every later block. A model can fail either criterion, and both outcomes remain in the table. The sixty-four-block design extends the original prospective window without treating the absence of a later reversal as evidence about unobserved future blocks. Reuse adoption must combine these measured costs with the policy-cost interval and the decision maker's explicit computation price and replacement fee.

\subsection{Two controls and horizon dependence}
\input{tables/vector69}
The additional vector upper witness and the unchanged triangular lower cover implement Proposition~\ref{prop:vector-witness69}. Both learned arms use the same paid label states. The comparison distinguishes the effects of a proposal from the effects of uniform or adaptive action refinement. It is not a derivative branch-and-bound comparison, and vector label fitting is not relabeled as certificate-aware scalar training.

The horizon-four record includes memory, floating and rational query counts, training or transfer identity and complete process time. The previously observed large descendant-tree cost remains a warning about horizon and innovation dependence even when the state lattice is absent. These finite path diagnostics are not a simultaneous population policy-cost comparison for vector controllers. The independent rational implementation checks every represented new mid-bin trajectory and independent scalar and two-control terminal integrals; it does not independently re-integrate every node of the largest inherited simplicial tree.

The resulting economic distinction is precise. The method returns an implementable policy under an original-optimum contract; an achieved readout certificate measures what the fitting procedure actually produced; a prospective work account prices its use; and centered evaluation can assess policy differences at a prespecified scale. None of those objects is substituted for the others. The original NBO construction and its broader applications retain their stated hypotheses, while the new numerical conclusions concern the particular expected-cost economies evaluated here.
